\section{Development kits interoperability and language boundaries}\label{sec:C22:adks}
An agent development kit typically bundles authoring conventions, tool integration, orchestration, testing, and deployment adapters. This breadth can improve consistency but also creates version coupling. A task object serialized by one language implementation may not preserve every field expected by another. A model adapter can normalize errors differently, changing retry behavior. Cross-language support is therefore a compatibility question rather than a guarantee that all implementations have identical semantics.

Google ADK's inspected repository describes graph-based workflows and structured delegation surfaces, while Microsoft Agent Framework describes multi-language agent and workflow construction~\cite{executionproduct003,executionproduct007}. These are concrete interfaces for conformance tests. The experiment should pin the relevant implementation and compare equivalent task graphs, state updates, and failure injections. It should not compare unrelated examples and attribute the difference to the toolkit brand.

Protocol support is similarly conditional. A framework advertising MCP can expose tools, resources, or prompts through different transport and authorization configurations. An A2A integration can exchange task status without sharing the same memory model or trust boundary. Chapter~\ref{ch:C19} treats these protocol distinctions directly. At the framework layer, the important question is which semantics are preserved when an adapter maps between local objects and external messages.

Serialization requires particular care for authority-bearing data. A human approval record should identify the precise operation, relevant arguments, principal, expiry, and policy epoch. Mapping it to an unstructured string in another framework weakens its meaning. The receiving application must validate the approval against its own action boundary. Successful message delivery cannot establish that the transferred object authorizes the requested effect.

Interoperability experiments should include a deliberate version mismatch, optional fields unknown to the receiver, partial tool failures, and revoked credentials. Report which failures are detected before dispatch and which emerge only after an effect. A graceful error is often preferable to a silent fallback that changes task meaning. The strongest interoperability result is a documented set of preserved invariants under declared adapters, not an assertion that any agent can collaborate safely with any other.

